\begin{table}[!htbp]
\centering
\caption{Effects evaluated at observed exposure support under alternative 3 km mappings}
\label{tab:ulez_boundary_effects}
\small
\begin{threeparttable}
\begin{tabularx}{0.96\textwidth}{>{\raggedright\arraybackslash}X >{\centering\arraybackslash}m{3.4cm} >{\centering\arraybackslash}m{3.4cm}}
\toprule
Estimand & Station-based & Boundary-based \\
\midrule
Static total treated effect & \shortstack{\(-9.22\%\)\\{\scriptsize \([-12.97,-5.30]\)}} & \shortstack{\(-8.56\%\)\\{\scriptsize \([-12.53,-4.40]\)}} \\
Dynamic entrants, not previously exposed & \shortstack{\(-4.58\%\)\\{\scriptsize \([-7.37,-1.71]\)}} & \shortstack{\(-4.54\%\)\\{\scriptsize \([-7.34,-1.66]\)}} \\
Dynamic entrants, previously exposed & \shortstack{\(-6.83\%\)\\{\scriptsize \([-11.91,-1.47]\)}} & \shortstack{\(-4.96\%\)\\{\scriptsize \([-8.76,-1.01]\)}} \\
Dynamic currently untreated exposure & \shortstack{\(-6.70\%\)\\{\scriptsize \([-10.42,-2.83]\)}} & \shortstack{\(-3.99\%\)\\{\scriptsize \([-5.63,-2.32]\)}} \\
\bottomrule
\end{tabularx}
\begin{tablenotes}
\footnotesize
\item \emph{Notes:} Entries are implied percentage effects with 95 percent confidence intervals in brackets. Effects are evaluated at the observed mean exposure relevant to each estimand. Because the two mappings generate different exposure distributions, the magnitudes are mapping-specific and should not be interpreted as estimates of a single invariant parameter.
\end{tablenotes}
\end{threeparttable}
\end{table}
